% Ordinary of Missa XVII: \input by missa-xvii.tex and ../kyriale/kyriale.tex.
% Score paths are relative to a sibling directory, so they work from both.
% ---------------------------------------------------------------------------
\chant{Kyrie}{tonus I}
\lilypondfile{../missa-xvii/scores/kyrie.ly}
\begin{tekst}
\w{Kýrie, eléison. Christe, eléison. Kýrie, eléison.}
  {Panie, zmiłuj się nad nami. Chryste, zmiłuj się nad nami.
   Panie, zmiłuj się nad nami.}
\end{tekst}

{\centering\small\rubr{W Adwencie i Wielkim Poście nie śpiewa się Gloria.}\par}

\credo{i}

\clearpage
% ---------------------------------------------------------------------------
\chant{Sanctus}{tonus V}
\lilypondfile{../missa-xvii/scores/sanctus.ly}
\begin{tekst}
\w{Sanctus, Sanctus, Sanctus Dóminus Deus Sábaoth.}
  {Święty, Święty, Święty, Pan Bóg Zastępów.}
\w{Pleni sunt caeli et terra glória tua.}
  {Pełne są niebiosa i ziemia chwały Twojej.}
\w{Hosánna in excélsis.}{Hosanna na wysokości.}
\w{Benedíctus qui venit in nómine Dómini.}
  {Błogosławiony, który idzie w imię Pańskie.}
\w{Hosánna in excélsis.}{Hosanna na wysokości.}
\end{tekst}

% ---------------------------------------------------------------------------
\chant{Agnus Dei}{tonus V}
\lilypondfile{../missa-xvii/scores/agnus.ly}
\begin{tekst}
\w{Agnus Dei, qui tollis peccáta mundi: miserére nobis.}
  {Baranku Boży, który gładzisz grzechy świata, zmiłuj się nad nami.}
\w{Agnus Dei, qui tollis peccáta mundi: miserére nobis.}
  {Baranku Boży, który gładzisz grzechy świata, zmiłuj się nad nami.}
\w{Agnus Dei, qui tollis peccáta mundi: dona nobis pacem.}
  {Baranku Boży, który gładzisz grzechy świata, obdarz nas pokojem.}
\end{tekst}

% ---------------------------------------------------------------------------
\chant{Ite, missa est}{tonus IV}
\lilypondfile{../missa-xvii/scores/ite.ly}
\begin{tekst}
\w{\V Ite, missa est.}{\V Idźcie w pokoju Chrystusa.}
\w{\R Deo grátias.}{\R Bogu niech będą dzięki.}
\end{tekst}
